\section{Interpretation and limitations}

The sharp expected-length characterization in \cref{obj:thm:full-honest-frontier-answer} separates two sources of uncertainty in the homogeneous causal log odds. The numerator contribution depends on the combined smoothness of the propensity and prognosis logits, whereas the denominator contribution depends on prognosis smoothness and is multiplied by the evaluation radius. Through \cref{obj:def:diagnostic-envelope}, these contributions give
\[
d_n(\alpha,\beta;r)
=
n^{-\min\{1/2,\,2(\alpha+\beta)/(2\alpha+2\beta+1)\}}
+
r\,n^{-4\beta/(4\beta+1)}.
\]
The sharp comparison gives this decomposition an optimality interpretation under its stated conditions: it describes the smallest worst-case expected length on each radius slice among procedures satisfying ambient honesty. The construction in \cref{obj:thm:finite-honest-upper} attains the corresponding upper bound with a single procedure whose tuning uses the sample size and public smoothness exponents.

The distinction between coverage and length evaluation is essential to this interpretation. In \cref{obj:def:length-objective}, the evaluation radius restricts the laws over which expected length is assessed, while admissibility continues to require global \(90\%\) coverage over the ambient model in \cref{obj:def:honest-procedures}. Thus, as the evaluation radius approaches zero, the procedure remains accountable for coverage throughout that model. At radius zero, the denominator contribution vanishes and the numerator contribution remains. The sharp comparison therefore quantifies the precision available at a zero effect while retaining the same coverage obligation for nearby and larger effects.

This behavior follows the structure of the observable odds identity in \cref{obj:lem:observable-odds}. The numerator is proportional to the denominator times the exponential effect increment. Consequently, denominator uncertainty has less influence on the recovered log odds when the effect is small, while numerator uncertainty continues to enter the inversion. The separate coordinate estimates and ranks in \cref{obj:def:coordinate-estimates,obj:def:resolutions} accommodate these different roles. Joint inversion then converts the coordinate bounds into an interval whose expected length responds to effect magnitude.

The balance can be expressed using the numerator exponent \(\mathsf a(\alpha,\beta)\) and the denominator exponent \(\mathsf b(\beta)\) introduced in \cref{obj:def:diagnostic-envelope}. The two contributions are comparable when
\[
r\asymp n^{-\{\mathsf a(\alpha,\beta)-\mathsf b(\beta)\}}.
\]
When \(\mathsf b(\beta)<\mathsf a(\alpha,\beta)\), this scale decreases with sample size. Above it, the effect-scaled denominator contribution is larger; below it, the numerator contribution is larger. Rough prognosis can therefore govern expected length at appreciable effect magnitudes even when combined propensity and prognosis smoothness gives a faster numerator rate. Concentrating evaluation near zero progressively reduces that contribution.

The known uniform covariate design and public exponents make this experiment precise. The design in \cref{obj:ass:uniform-design} supplies the integration measure for the projection construction, fixing the geometry used to control approximation error. The public exponents specify the smoothness restrictions in \cref{obj:def:model} and determine the two resolutions before sampling. Together with the prescribed envelopes, these features support the finite-sample coverage guarantee and the comparison of expected lengths across effect radii. The rational realization in \cref{obj:lem:fixed-budget-realization} also gives a concrete numerical meaning to the construction: prescribed precision budgets produce outward endpoint enclosures with the stated additional-length guarantee.

\subsection{Limitations and future work}

The precision characterization concerns a known uniform design, public smoothness exponents, a scalar covariate, and a homogeneous conditional log-odds coefficient. Unknown design would require accounting for the covariate distribution in the projection construction. Unknown smoothness raises a separate adaptation problem, including the coverage requirements appropriate to a collection of smoothness classes. Higher-dimensional covariates would change the approximation problem and its balance with sampling uncertainty.

Heterogeneous effects would require specifying a new target, such as an average contrast or an effect function, together with an appropriate coverage criterion. The scalar inversion and radius-based length objective studied here provide a starting point for formulating such questions; their extensions require separate precision characterizations.

Practical calibration and computational optimization also merit further study. The proved construction supplies coverage and expected-length guarantees with prescribed arithmetic precision. Empirical assessment could examine interval lengths at realistic sample sizes, sensitivity to conservative constants, and numerical performance across admissible implementations. Software development could optimize evaluation of the projection statistics and endpoint enclosures while preserving the stated guarantees. These prospective assessments concern practical performance beyond the proved precision characterization.
